% Generated draft: requires mathematical and source review.
\documentclass[11pt]{article}
\usepackage[margin=25mm]{geometry}
\usepackage{fontspec}
\setmainfont{Latin Modern Roman}
\usepackage{amsmath,amssymb,mathtools,unicode-math}
\setmathfont{Latin Modern Math}
\usepackage{enumitem,needspace,xurl,hyperref}
\setlength{\parindent}{0pt}
\setlength{\parskip}{5pt}
\setlength{\emergencystretch}{3em}
\newcommand{\N}{\mathbb{N}}
\newcommand{\Z}{\mathbb{Z}}
\newcommand{\Q}{\mathbb{Q}}
\newcommand{\R}{\mathbb{R}}
\newcommand{\C}{\mathbb{C}}
\newcommand{\F}{\mathbb{F}}
\newcommand{\A}{\mathbb{A}}
\newcommand{\PP}{\mathbb{P}}
\newcommand{\E}{\mathbb{E}}
\newcommand{\eps}{\varepsilon}
\newcommand{\dd}{\,\mathrm{d}}
\begin{document}
\section*{UnsolvedMath ID: 30005340}
\textbf{Source problem code:} catalogue entry
\subsection{Definitions and mathematical statement}
The supplied record poses a research programme rather than a fully specified conjecture. The following notation makes its unresolved choices explicit.

For a specified random function \(f\) on a specified set \(D\subseteq\R\), define the distinct-zero count by
\[
N_D^{\mathrm{dist}}(f)=\begin{cases}
\#\{x\in D:f(x)=0\},&\text{if this set is finite},\\
+\infty,&\text{otherwise}.
\end{cases}
\]
This definition includes the identically zero function. Whenever this extended nonnegative random variable is measurable, its expectation is defined in \([0,+\infty]\). For a nonzero polynomial \(P\), an alternative convention is
\[
N_D^{\mathrm{mult}}(P)=\sum_{\substack{x\in D\\P(x)=0}}\operatorname{mult}_x(P),
\]
where \(\operatorname{mult}_x(P)\) is the algebraic multiplicity of the root. The supplied record does not specify which convention its notation \(N_{n,\xi}\) uses, or how a positive-probability identically zero polynomial is treated.

\textbf{Requested mathematical target.} For specified ensembles of random polynomials indexed by \(n\), with coefficient randomness denoted by \(\xi\), determine the expected number
\[
\E N_{n,\xi}
\]
of real roots beyond the Gaussian coefficient case; more generally, determine the expected number of real zeros for specified ensembles of random functions. The source does not specify whether the desired determination is an exact formula, an asymptotic equivalent as \(n\to\infty\), or an estimate with a prescribed error.

\textbf{Unresolved qualifications.} The complete supplied excerpt does not define the polynomial ensemble, the deterministic coefficient weights, the coefficient laws or dependence assumptions, the meaning of the index \(n\), or the random-function class and its domain. It also supplies no regularity, measurability, nondegeneracy, or integrability hypotheses ensuring a well-defined finite expected zero count. Consequently, the notation above is a conditional framework for the requested target, not a claim that these missing choices have been recovered. The mean-zero and unit-variance assumptions appearing in the historical discussion, and the independence assumptions in its central-limit-theorem illustration, cannot be imposed on the entire problem from this excerpt alone.

The excerpt reports Gaussian-case formulas and universality comparisons, including a displayed comparison of the form
\[
\E N_{n,\xi}=(1+o(1))\E N_{n,\mathrm{Gauss}}.
\]
Its surrounding text is interleaved and incomplete, so the ensemble and hypotheses governing this comparison cannot be reliably assigned. It is supplied historical context, not a universally quantified assertion or an additional conjecture here.
\subsection{Short English statement}
Determine the expected number of real roots of specified non-Gaussian random polynomial ensembles, and seek analogous results for random functions. The supplied source leaves the ensembles, coefficient hypotheses, root-counting convention, function domains, and intended precision unspecified, so it does not determine a single precise quantitative conjecture. Its \texttt{partially\_solved} label is supplied dated curation requiring independent review.
\subsection{Sources}
\begin{itemize}
\item Local registry: \texttt{data/unsolvedmath/problems.json}, with \texttt{catalog/problems.json} as the supplementary registry; supplied permanent UnsolvedMath ID \(30005340\), title \emph{Expected Real Roots of General Random Polynomials}, source code \texttt{catalogue entry}. The supplied excerpt begins with the label ``Problem 3.'' No individual proposer is identified.
\item Supplied report citation: \emph{Combinatorics} (2023), Oberwolfach Reports, DOI \texttt{10.4171/owr/2023/1}, pp.~5--89. \url{https://doi.org/10.4171/owr/2023/1}. This is also the supplied literature-source link labelled \emph{Oberwolfach Reports 2023/1}. No problem-specific page is supplied.
\item Supplied literature reference: Do--Nguyen, \emph{Real Roots of Random Polynomials}. \url{https://arxiv.org/abs/2303.05478}. The author attribution and title are reproduced as supplied; their correspondence to the linked document has not been independently verified.
\item Literature assessment dated exactly \texttt{2026-08-22}: the supplied curation labels the problem partially solved, reports expectation formulas for Kac, Weyl, and broad coefficient classes, and states that no theorem determines the expectation for general random polynomials or arbitrary random functions. This is supplied dated curation, not an independently verified literature assessment. The background expressly requests independent verification of time-sensitive status and cited claims.
\item The record labels its statement \texttt{corrected\_verified} and reports a check against the official Oberwolfach report. That verification claim is part of the supplied record and has not been independently checked here.
\end{itemize}
\end{document}
